\chapter{LUT, FF, Slice와 전용 자원의 기본기}

\begin{summarybox}[이 장의 목표]
Verilog 연산자를 보기 전에 그 연산을 받아 주는 FPGA 자원을 먼저 이해한다.
LUT는 조합논리의 1-bit 출력을 만들고, FF는 1-bit 상태를 저장하며, slice는
이들을 carry와 local mux와 함께 배치하는 물리 단위다. DSP와 BRAM은 많은
LUT를 대신하는 hard block이며, LUTRAM과 SRL은 LUT를 논리 이외의 용도로
사용한다.
\end{summarybox}

\section{합성은 코드를 실행 순서로 번역하지 않는다}

합성기는 RTL을 한 줄씩 실행하는 프로그램으로 만들지 않는다. 조합식은 입력이
바뀔 때마다 값을 만드는 회로가 되고, clocked assignment는 값을 저장하는
register가 된다. 그 사이에 있는 Boolean function, arithmetic, selection과
memory access를 목표 FPGA가 가진 primitive에 맞게 mapping한다.

\begin{lstlisting}
wire [22:0] sum_w = a + b;
reg  [22:0] sum_q;

always @(posedge clk)
  if (en)
    sum_q <= sum_w;
\end{lstlisting}

이 예에서 \rtlfile{sum\_w}는 별도의 저장 공간이 아니다. \rtlfile{a+b}를
구현하는 조합 회로의 net 이름이다. \rtlfile{sum\_q}는 clock edge 사이에서
값을 보존해야 하므로 23개의 FF가 필요하다. \rtlfile{en}은 clock 자체를
논리로 만들기보다 FF의 clock enable로 mapping될 수 있다.

\section{LUT6: 1-bit 출력 함수를 저장하는 진리표}

7-series FPGA의 LUT6는 최대 여섯 입력에 대한 1-bit Boolean function을
구현한다. 내부 동작은 AND와 OR gate를 RTL 순서대로 연결하는 것보다, 입력을
주소로 사용해 미리 정해 둔 1-bit 결과를 고르는 작은 truth-table memory로
생각하는 편이 정확하다.

\begin{lstlisting}
assign y = (a & b) | (~c & d);
\end{lstlisting}

위 식은 네 입력과 한 출력의 함수이므로 주변 논리와 합쳐지지 않는다는 가정에서
LUT6 하나에 들어갈 수 있다. 식에 연산 기호가 세 개 있다는 이유로 LUT 세 개를
세면 안 된다. 합성기는 전체 Boolean function을 최소화하고 LUT 입력 수에 맞춰
packing한다.

반대로 다음 23-bit XOR는 LUT 하나가 아니다.

\begin{lstlisting}
assign y[22:0] = a[22:0] ^ b[22:0];
\end{lstlisting}

각 \rtlfile{y[i]}는 독립된 출력이다. LUT6의 기본 출력은 한 bit이므로 구조적
상한을 잡을 때에는 1-bit XOR가 폭만큼 병렬 복제된다고 본다. 다른 논리와
흡수되지 않은 isolated 23-bit XOR라면 최대 23 LUT 수준에서 시작한다.

\begin{summarybox}[LUT를 셀 때의 두 질문]
\begin{enumerate}
  \item 각 output bit를 만드는 데 서로 다른 입력 변수가 몇 개 필요한가?
  \item 서로 독립된 output bit가 몇 개 필요한가?
\end{enumerate}
첫 질문은 LUT level과 분할 가능성을, 두 번째 질문은 대략적인 LUT 폭을 알려 준다.
\end{summarybox}

하나의 LUT는 입력을 공유하는 두 개의 작은 함수로 분할되어 LUT6\_2 형태로
사용될 수도 있다. 따라서 RTL에서 센 Boolean output 수와 최종 LUT 수가 항상
같지는 않다. common subexpression, constant propagation, downstream trimming,
dual-output packing이 최종 수를 바꾼다. 이 때문에 면적 추정은 범위를 잡는
작업이고, 정확한 값은 hierarchy utilization에서 확인해야 한다.

내가 LUT를 셀 때에는 RTL의 연산 기호부터 세지 않는다. 먼저 output bit 수를
세고, 각 output에 들어오는 독립 입력 수를 적는다. 그다음 carry, MUXF, DSP 또는
memory로 빠질 구조를 제외한다. 이 순서로 계산한 값은 정확한 report 수가 아니라
합성 결과가 상식적인 범위인지 판단하는 기준이 된다.

\section{FF와 FDRE: 1-bit 상태와 control set}

FF는 clock edge에서 1 bit를 저장한다. 다음 register는 논리적으로 23 FF다.

\begin{lstlisting}
reg [22:0] data_q;
always @(posedge clk) begin
  if (rst)
    data_q <= 23'd0;
  else if (ce)
    data_q <= data_d;
end
\end{lstlisting}

7-series에서 이러한 동기 reset과 clock enable을 가진 FF는 FDRE 계열 자원으로
mapping될 수 있다. 그러나 FF 개수와 slice 개수는 같은 면적 지표가 아니다.
같은 clock, reset, enable을 공유하는 register들은 같은 control set을 가져
가까운 slice에 pack되기 쉽다. 같은 수의 FF라도 reset polarity나 enable이
여러 종류로 갈라지면 packing 후보가 제한된다.

\begin{lstlisting}
// One shared control set
always @(posedge clk) begin
  if (rst) begin
    a_q <= 0;
    b_q <= 0;
  end else if (ce) begin
    a_q <= a_d;
    b_q <= b_d;
  end
end
\end{lstlisting}

모든 register에서 reset을 무조건 제거하거나 하나로 합치는 것이 목표는 아니다.
protocol state처럼 reset 값이 기능에 필요한 bit는 유지해야 한다. 반면 data
pipeline처럼 valid가 0일 때 값이 관찰되지 않는 register는 data reset이 없어도
되는지 검토할 수 있다. 이 결정은 면적뿐 아니라 SRL/DSP/BRAM 내부 register
inference에도 영향을 준다.

내 설계에서 FF를 줄일 때에도 register 선언 개수만 지우지 않았다. pipeline의
각 cycle에서 실제로 살아 있는 data, valid, mode와 address bit를 표로 만든 뒤,
그 stage 이후에 사용되지 않는 rail만 제거했다. 그래야 면적은 줄고 cycle
alignment는 유지된다.

\begin{warningbox}[FF가 공짜라는 오해]
FPGA에 LUT보다 FF가 남는 경우가 많아도 FF를 멀리 떨어진 stage마다 추가하면
clock, reset, enable 배선과 placement가 함께 변한다. pipeline register는
critical path를 끊는 위치와 data/valid/address 정렬을 정의할 때 가치가 있다.
단순히 utilization 비율이 낮다는 이유로 넓은 복제 register를 늘리지 않는다.
\end{warningbox}

\section{Slice, SLICEL과 SLICEM}

Slice는 LUT, FF, carry logic과 local mux를 함께 배치하는 물리적 묶음이다.
SLICEL은 일반 logic과 register를 제공하고, SLICEM은 그 기능에 더해 일부 LUT를
distributed RAM이나 SRL로 사용할 수 있다. 그래서 다음 두 설계는 LUT 수가
비슷해도 placement 제약이 다를 수 있다.

\begin{itemize}
  \item 넓은 Boolean datapath는 SLICEL과 SLICEM의 logic LUT를 사용할 수 있다.
  \item LUTRAM이나 SRL이 많은 설계는 SLICEM 위치를 필요로 한다.
  \item 긴 carry chain은 인접 slice의 전용 수직 연결을 따라 배치되는 것이 좋다.
  \item BRAM과 DSP 주변 register는 hard block column과의 거리가 timing에 영향을 준다.
\end{itemize}

따라서 post-synthesis LUT/FF 수만으로 route 결과를 예측할 수 없다. 같은 논리
개수라도 control set, SLICEM 수요, carry/DSP/BRAM column, congestion이 다르면
post-route WNS가 달라진다.

\section{CARRY4: 덧셈과 뺄셈의 전용 전달 경로}

일반 LUT와 routing만으로 각 bit의 carry를 다음 bit로 전달하면 넓은 산술 회로의
지연이 커진다. CARRY4는 네 bit의 carry를 slice 내부 전용 경로로 전달한다.
LUT는 각 bit의 propagate/select 성격을 만드는 데 쓰이고, carry primitive가
carry와 sum 생성을 담당한다.

\begin{lstlisting}
wire [22:0] add_y = a + b;
wire [22:0] sub_y = a - b;
\end{lstlisting}

23-bit add/subtract는 구조적으로 $\lceil23/4\rceil=6$개의 CARRY4 구간을
생각할 수 있다. 이번 isolated synthesis에서도 add와 subtract가 각각 LUT 23,
CARRY4 6으로 나왔다. 이것은 고정 법칙이 아니라 이 device와 합성 문맥에서 나온
측정값이다. 상수 operand, 사용되지 않는 MSB, 주변 logic 흡수에 따라 달라진다.

이 수치를 얻고 난 뒤에는 23-bit add/subtract를 보면 먼저 CARRY4 여섯 구간을
떠올린다. 두 번의 correction이 한 cycle에 있으면 열두 구간 수준까지 이어질 수
있다고 예상하고 timing report에서 실제 carry primitive 수를 확인한다.

폭 하나만으로 timing을 판단해서도 안 된다. 다음 경로들은 같은 23-bit adder를
포함해도 서로 다르다.

\begin{itemize}
  \item register $\rightarrow$ adder $\rightarrow$ register
  \item register $\rightarrow$ mode MUX $\rightarrow$ adder $\rightarrow$ register
  \item register $\rightarrow$ adder $\rightarrow$ comparator $\rightarrow$ correction MUX
        $\rightarrow$ register
\end{itemize}

Unified NTT의 150~MHz 실패 사례도 단순 subtract 하나가 아니라 carry chain
두 개와 mode-dependent selection이 같은 register 경계에 들어간 경로였다.

\section{MUXF7/MUXF8과 넓은 선택기}

한 output bit의 2:1 MUX는 두 data와 select, 즉 세 입력 함수다. 4:1 MUX도
네 data와 두 select의 여섯 입력 함수이므로 한 LUT6에 들어갈 수 있다.

\begin{lstlisting}
assign y2 = sel1 ? d1 : d0;
assign y4 = (sel2 == 2'd0) ? d0 :
            (sel2 == 2'd1) ? d1 :
            (sel2 == 2'd2) ? d2 : d3;
\end{lstlisting}

실제 isolated synthesis에서 23-bit 2:1과 4:1 MUX가 모두 LUT 23이었다. 각
output bit의 함수가 LUT6 하나에 들어갔기 때문이다. 입력 수가 더 늘어나면 여러
LUT output을 MUXF7/MUXF8 같은 slice-local mux로 합칠 수 있고, 그 범위를 넘거나
배치가 맞지 않으면 추가 LUT level과 general routing이 생긴다.

MUX 면적은 대략 ``output 폭 $\times$ 필요한 선택 단계''로 시작해 생각한다.
그러나 timing은 select fanout, data source 위치, arithmetic 앞뒤 배치까지 본다.
23 LUT라는 같은 utilization 결과가 같은 delay를 의미하지 않는다.

처음에는 4:1 MUX가 2:1 MUX보다 LUT를 두 배쯤 쓸 것으로 예상했지만, 23-bit
실험에서는 둘 다 23 LUT였다. 네 data와 두 select가 정확히 LUT6 입력 여섯 개에
들어갔기 때문이다. 대신 select fanout과 input pin, 배치가 다르므로 timing까지
같다고 판단하지 않는다.

\section{DSP48E1: 곱셈기와 내부 pipeline}

DSP48E1은 단순 multiplier만이 아니라 operand register, multiplier, ALU와 output
register를 가진 hard block이다. 다음과 같이 산술 의도를 직접 쓰면 합성기가
operand 폭과 사용 가능한 DSP 구조를 보고 mapping할 수 있다.

\begin{lstlisting}
(* use_dsp = "yes" *) wire [38:0] product = a23 * b16;
\end{lstlisting}

이번 23x16 isolated experiment는 fabric LUT 0, DSP48E1 1개로 합성되었다.
입출력 register 일부도 DSP 내부로 흡수될 수 있어 top-level FF 수만 보고 pipeline
위치를 단정하면 안 된다. synthesis schematic과 DSP property, timing path의
primitive pin을 함께 확인한다.

이 실험에서 23$\times$16 곱셈은 LUT를 쓰지 않고 DSP48E1 하나가 됐다. 그래서
곱셈 기호를 보면 무조건 비싸다고 판단하지 않는다. operand 폭이 DSP에 맞는지,
DSP budget이 남는지, 내부 register가 어느 stage에 흡수됐는지를 먼저 본다.

직접 shift-add 식으로 풀면 DSP 사용을 피할 수 있지만 adder tree와 register가
fabric에 생긴다. DSP가 부족한 설계에서는 유효한 교환이고, DSP가 남는 설계에서는
LUT와 routing을 낭비할 수 있다. Unified NTT는 네 개 DSP를 유지하면서 여러
scheme의 multiplication을 stage와 mode로 공유한다.

\section{RAMB36/RAMB18, LUTRAM과 SRL}

큰 coefficient memory와 ROM은 RAMB36E1/RAMB18E1 같은 block RAM에 들어가는 것이
일반적이다. 작은 비동기 read memory는 LUTRAM이 적합할 수 있고, 단순 delay line은
SRL로 들어갈 수 있다.

\begin{longtable}{L{30mm}L{50mm}L{78mm}}
\toprule
자원 & 잘 맞는 구조 & 확인할 조건 \\
\midrule
\endhead
BRAM & 깊은 RAM/ROM, FIFO & synchronous read, port 수, width/depth, write mode \\
LUTRAM & 작은 distributed RAM & SLICEM 수요, asynchronous/synchronous read 형태 \\
SRL & reset 없는 delay line & tap 위치, enable, 개별 stage access 여부 \\
FF chain & stage마다 상태가 필요한 pipeline & reset/enable, 각 tap 사용, control alignment \\
\bottomrule
\end{longtable}

예를 들어 23-bit 값을 세 cycle 미루는 구조는 논리적으로 69 FF다. 중간 tap을
쓰지 않고 data reset도 필요 없다면 SRL 후보가 될 수 있다. 반대로 각 stage가
서로 다른 valid나 mode와 결합되거나 중간 값을 butterfly가 사용하면 독립 FF
pipeline이 필요하다.

Unified NTT에서 coefficient memory는 계산에 필요한 동시 접근 수와 bank conflict를
기준으로 BRAM bank를 구성한다. 최종 standalone의 6 BRAM tile은 단순히 bit 수를
저장한 결과가 아니라 port와 병렬 lane 요구를 포함한 구조다. Stream complete
system의 추가 3 tile은 accelerator memory가 아니라 DMA payload FIFO에서 나온다.

작은 합성 예제에서도 32$\times$8 asynchronous-read memory는 LUTRAM 8개,
576$\times$23 synchronous memory는 RAMB36E1 하나가 됐다. 내가 BRAM 면적을
total bit만으로 나누지 않고 width, depth와 port를 함께 적는 이유다.

\section{자원 수를 회로 그림으로 바꾸는 순서}

\begin{checklistbox}
\begin{enumerate}
  \item 조합 출력 bit 수와 각 bit의 독립 입력 수를 세어 LUT 폭과 level을 추정한다.
  \item 저장해야 하는 live bit와 control set을 세어 FF와 packing을 추정한다.
  \item add/subtract/magnitude compare가 만드는 carry chain 길이를 본다.
  \item 넓은 selection이 LUT6 한 단계인지, MUXF 계층이나 추가 LUT가 필요한지 본다.
  \item multiply와 multiply-add가 DSP 폭과 내부 pipeline에 맞는지 확인한다.
  \item memory의 depth만 보지 말고 width, port, read latency를 함께 적는다.
  \item delay line에 reset과 중간 tap이 있는지 보고 SRL과 FF chain을 구분한다.
  \item 마지막으로 synthesis hierarchy와 post-route timing에서 이 추정을 검증한다.
\end{enumerate}
\end{checklistbox}
